\chapter{Évaluer et mesurer}
\label{ch:evaluation}

La qualité technique vaut 20\,\% et s'intitule \og résultats mesurés et bien évalués \fg ; la
frugalité vaut 20\,\% et se mesure. Ce chapitre fixe \emph{avant} de coder ce qu'on mesure, sur
quoi, et comment. Les tableaux sont des gabarits : ils seront remplis mardi soir et présentés tels
quels.

\section{Ce qu'on mesure}

\begin{center}\small
\begin{tabularx}{\textwidth}{@{}L{3.6cm} Y L{3.2cm} l@{}}
\toprule
\textbf{Mesure} & \textbf{Comment} & \textbf{Sur quoi} & \textbf{Critère} \\
\midrule
Précision, F1 macro & scikit-learn, par intention et par famille & Tests Banking77 et nos tests &
  Technique \\
Taux de réussite & Rejeu automatique des scénarios & 30 conversations & Technique, métier \\
Échecs du garde-fou & Compteur dans la trace & Toutes les réponses générées & Explicabilité \\
Répartition par niveau & Compteur dans la trace : mots-clés, classifieur, escalade & Tests et
  scénarios & Frugalité \\
Taux d'escalade & Compteur dans la trace & Scénarios & Frugalité \\
Part sans LLM & Compteur dans la trace & Scénarios & Frugalité \\
Latence & Horloge par étage (p50, p95) & Scénarios & Frugalité \\
Énergie & CodeCarbon, par lot (section~\ref{sec:energie}) & Lots de 200 requêtes & Frugalité \\
Taille & Go sur disque, quantisation comprise & Chaque modèle chargé & Frugalité \\
Appétence & Aire sous la courbe ROC & 20\,\% de Bank Marketing tenus à part & Technique \\
\bottomrule
\end{tabularx}
\end{center}

\section{La compréhension}

\paragraph{Trois jeux de test.} (a) Le test officiel de Banking77 en anglais, 3\,080 phrases ;
(b) le même, traduit en français ; (c) nos 160 phrases françaises des 8 nouvelles intentions,
écrites à la main. Le jeu (a) rend nos chiffres comparables à l'article d'origine ; les jeux (b) et
(c) mesurent ce qui compte pour le sujet.

\begin{center}\small
\begin{tabularx}{\textwidth}{@{}Y c c c c c@{}}
\toprule
\textbf{Configuration} & \textbf{Préc. (a)} & \textbf{Préc. (b)} & \textbf{F1 famille (b+c)} &
  \textbf{Latence p95} & \textbf{Taille} \\
\midrule
\textit{Référence : BERT affiné (article)} & \textit{93,66\,\%} & --- & --- & --- & \textit{1,4 Go} \\
Mots-clés seuls, sur la part acceptée & \amesurer & \amesurer & \amesurer & \amesurer & Ko \\
TF-IDF + régression logistique & \amesurer & \amesurer & \amesurer & \amesurer & \amesurer \\
MiniLM-L12 + régression logistique & \amesurer & \amesurer & \amesurer & \amesurer & 0,47 Go \\
\textbf{e5-small + régression logistique} & \amesurer & \amesurer & \amesurer & \amesurer &
  \textbf{0,47 Go} \\
e5-base + régression logistique & \amesurer & \amesurer & \amesurer & \amesurer & 1,1 Go \\
\f{qwen3:1.7b}, sans exemple & \amesurer & \amesurer & \amesurer & \amesurer & 1,4 Go \\
\f{ministral-3:14b}, sans exemple & \amesurer & \amesurer & \amesurer & \amesurer & 9,1 Go \\
\bottomrule
\end{tabularx}
\end{center}

Pour les mots-clés, on donne aussi la part des messages qu'ils acceptent : l'organisateur l'estime
à environ 30\,\%, nous la mesurons. La taille de la référence correspond à BERT-Large (340 millions de paramètres en float32). Les
deux dernières lignes classent avec un LLM, la liste des intentions dans la consigne : elles
montrent ce que coûterait la compréhension \og tout LLM \fg. Le message attendu : un encodeur de
118 millions de paramètres fait aussi bien, ou presque, pour une fraction du coût.

\paragraph{Le seuil $\tau$.} Sur un jeu de validation (10\,\% de l'entraînement, tenus à part), on
trace la précision des messages acceptés en fonction de la part de messages acceptés, pour $\tau$
de 0,30 à 0,95. On retient le plus petit $\tau$ qui donne au moins 97\,\% de précision. La courbe
elle-même va dans le support : elle montre qu'on a \emph{choisi} le compromis.

\section{Les scénarios de bout en bout}

Un scénario est une conversation complète, avec un client donné et des attentes vérifiables
automatiquement.

\begin{code}
\footnotesize
\begin{verbatim}
# eval/scenarios/sophie_epargne.yaml
client: C03
tours:
  - client: "je voudrais mettre de cote pour mes enfants"
    attendu: {famille: epargne, demande: [horizon]}
    interdit: {demande: [age, enfants]}        # deja connus par le profil
  - client: "dans 10 ans pour leurs etudes, 100 euros par mois"
    attendu: {propose: [livret_a], ecarte: [lep]}
    interdit: {nombre_invente: true}
\end{verbatim}
\end{code}

\paragraph{Le jeu.} Trente conversations : les neuf parcours, chacun en trois variantes
(formulation différente, autre client, réponse en plusieurs tours), plus trois demandes hors
périmètre. Chaque membre écrit les scénarios de son module ; Sohayb maintient le programme qui les
rejoue. Le jeu est rejoué à chaque intégration : \textbf{un scénario qui passait et qui échoue
bloque la fusion}.

\paragraph{Ce qui est vérifié.} La famille reconnue ; les questions posées, et surtout celles qui
\emph{ne doivent pas} l'être (une information déjà connue) ; les produits proposés et écartés ;
l'action déclenchée et sa confirmation ; le passage de main ; l'absence de nombre inventé.

\section{La formulation}

Deux indicateurs. \textbf{Le taux d'échec du garde-fou}, automatique, sur toutes les réponses
générées. \textbf{Une évaluation humaine courte} : pour chaque modèle candidat, du 14B au plus petit, 30 réponses notées
par deux membres sur trois critères (exact, clair, ton juste), en oui ou non, \textbf{champs
facultatifs} --- la leçon du formulaire de Timso AI.

\section{La frugalité : trois configurations}

Le cœur de la démonstration de frugalité est une comparaison sur le même jeu de scénarios.

\begin{center}\small
\begin{tabularx}{\textwidth}{@{}l Y@{}}
\toprule
\textbf{Configuration} & \textbf{Ce qu'elle fait} \\
\midrule
A --- Tout LLM & \f{ministral-3:14b} reçoit le message, le profil et le catalogue dans sa
  consigne, et fait tout : comprendre, décider, rédiger. C'est ce que ferait une équipe pressée. \\
B --- Cascade & Notre architecture complète. \\
C --- Cascade sans LLM & La même, avec des gabarits à la place du petit LLM : la borne basse en
  énergie. \\
\bottomrule
\end{tabularx}
\end{center}

\begin{center}\small
\begin{tabular}{@{}l c c c c c c@{}}
\toprule
& \textbf{Réussite} & \textbf{Nombres inventés} & \textbf{Wh / 1\,000 req.} &
  \textbf{Latence p95} & \textbf{Go chargés} & \textbf{Sans LLM} \\
\midrule
A & \amesurer & \amesurer & \amesurer & \amesurer & 9,1 & 0\,\% \\
B & \amesurer & \amesurer & \amesurer & \amesurer & 3,5 & \amesurer \\
C & \amesurer & \amesurer & \amesurer & \amesurer & 0,5 & 100\,\% \\
\bottomrule
\end{tabular}
\end{center}

Ce qu'on s'attend à montrer, et qu'il faudra confirmer : B atteint la qualité de A, voire mieux
parce que ses décisions sont des règles ; B consomme une fraction de l'énergie de A ; C coûte
encore moins mais rédige moins naturellement --- ce qui justifie, mesure à l'appui, la présence
d'un petit LLM. Si la mesure dit autre chose, c'est la mesure qu'on présente.

\section{Le protocole de mesure de l'énergie}
\label{sec:energie}

CodeCarbon mesure l'énergie de la machine entière (processeur, mémoire, GPU Nvidia par NVML). Une
requête dure trop peu pour être mesurée seule : on mesure par lots.

\begin{enumerate}
  \item Charger les modèles et faire une requête de chauffe (le chargement n'est pas compté).
  \item Mesurer 60 secondes à vide : puissance de repos $P_0$.
  \item Rejouer un lot de $N = 200$ requêtes tirées des scénarios ; énergie totale $E$, durée $t$.
  \item Énergie par requête : $(E - P_0 \times t) / N$.
  \item Répéter trois fois ; présenter la moyenne et l'écart.
\end{enumerate}

\begin{code}
\footnotesize
\begin{verbatim}
from codecarbon import OfflineEmissionsTracker
suivi = OfflineEmissionsTracker(country_iso_code="FRA", measure_power_secs=1,
                                project_name="cascade-B", log_level="error")
suivi.start()
for message, client in lot:
    reply(message, client, {})
suivi.stop()
energie_kwh = suivi.final_emissions_data.energy_consumed
\end{verbatim}
\end{code}

La version \og hors ligne \fg{} de CodeCarbon n'interroge aucun service extérieur : cohérent avec
la règle. On présente d'abord l'énergie en wattheures, puis l'équivalent CO\textsubscript{2}
calculé par CodeCarbon pour le réseau français. Sur Jetson, si le temps le permet, la puissance est
lue avec \f{tegrastats}.
